\documentclass{amsart}
\usepackage{amsmath,amssymb,amsthm,hyperref}

\newtheorem{theorem}{Theorem}
\newtheorem{proposition}{Proposition}
\newtheorem{lemma}{Lemma}
\newtheorem{corollary}{Corollary}
\theoremstyle{definition}
\newtheorem{definition}{Definition}
\newtheorem{assumption}{Assumption}
\theoremstyle{remark}
\newtheorem{remark}{Remark}

\begin{document}

\title{A Multilevel Mediation Model of Team Learning Behavior in Field Settings}
\maketitle

\begin{abstract}
We give an explicit, formally specified model of the antecedents of team
learning behavior in ongoing organizational work teams. The model
identifies factors at the individual, team, and organizational levels;
articulates the mechanism by which each acts; and singles out a shared
interpersonal-climate variable---team psychological safety, defined as the
shared belief that the team is safe for interpersonal risk taking---as the
mediator linking contextual supports (supportive leadership, contextual
support, coaching) to observed learning behavior. We state the model as a
recursive structural (path) system and prove, under explicitly stated
behavioral axioms and a strict-exogeneity (no-confounding) identification
condition, that (i) the structural coefficients are identified from
population regressions; (ii) the support--learning effect decomposes into a
direct part $\zeta$ and a strictly positive mediated part $\delta\beta$
running through psychological safety; (iii) the climate variable is a
\emph{complete} mediator if and only if the testable restriction $\zeta=0$
holds---a restriction motivated, but not logically forced, by the
mechanism lemmas---while partial mediation ($\zeta\neq0$) is the general
case; and (iv) the mediated chain raises team performance by
$\mu\delta\beta\gg0$. We are careful to distinguish what the behavioral
axioms entail (the signs $\delta>0$, $\beta\gg0$, $\mu>0$ and the existence
and positivity of the mediated path) from the separate empirical hypothesis
of complete mediation ($\zeta=0$).
\end{abstract}

\section{Formalization}

\subsection{Units and levels}
We consider a population $\mathcal{T}$ of ongoing work teams $t$, each
embedded in an organization and persisting over time, with member set
$M_t$, $|M_t|=n_t\ge 2$. Quantities indexed by $i$ range over individual
members $i\in M_t$; quantities indexed by $t$ are \emph{team-level}
(either properties of the team as a whole, e.g.\ leader behavior, or
aggregates of member-level quantities). This separation of levels is the
formal counterpart of the requirement to identify individual-, team-, and
organizational-level antecedents.

\subsection{The dependent construct: learning behavior}
\begin{definition}[Learning behavior]\label{def:LB}
For team $t$, let $B_t=(B_t^{1},\dots,B_t^{5})$ record the team's
\emph{learning behavior}, where the five coordinates are the (suitably
scaled, e.g.\ frequency-per-unit-time) intensities of the five activities
named in the problem:
\[
B_t^{1}=\text{asking questions},\quad
B_t^{2}=\text{seeking feedback},\quad
B_t^{3}=\text{discussing errors/unexpected outcomes},
\]
\[
B_t^{4}=\text{reflecting on past performance},\quad
B_t^{5}=\text{experimenting with new actions/methods}.
\]
Define the scalar \emph{learning-behavior level}
\[
L_t \;=\; \sum_{k=1}^{5} w_k\, B_t^{k},\qquad w_k>0,\ \sum_k w_k=1,
\]
a positively weighted index of the five activities. All five activities
are observable, so $L_t$ is an observable random variable.
\end{definition}

\begin{remark}
The defining feature shared by $B_t^1,\dots,B_t^5$ is that each requires a
member to expose a potential deficiency or to deviate from the routine in
front of others: a question reveals ignorance, feedback-seeking reveals
uncertainty about one's own performance, discussing an error reveals a
mistake, reflection surfaces collective shortcomings, and experimentation
risks visible failure. Hence every coordinate of $B_t$ is an
\emph{interpersonally risky} act. This single structural fact---proved as
Lemma~\ref{lem:risk} below from the definitions---is what makes an
interpersonal-climate variable the natural mediator.
\end{remark}

\subsection{The mediator: team psychological safety}
\begin{definition}[Psychological safety]\label{def:PS}
For member $i$ in team $t$ let $\psi_{t,i}\in\mathbb{R}$ denote $i$'s
belief about the consequences of taking interpersonal risks in the team,
measured so that larger values mean ``risk taking is safer'' (smaller
expected interpersonal cost). The team-level \emph{psychological safety}
is the shared component
\[
\Psi_t \;=\; \frac{1}{n_t}\sum_{i\in M_t}\psi_{t,i},
\]
and we say the belief is \emph{shared} when the within-team dispersion
$\operatorname{Var}_{i\in M_t}(\psi_{t,i})$ is small (formalized in
Assumption~\ref{ass:share}). $\Psi_t$ is the ``shared interpersonal
climate variable'' of the problem.
\end{definition}

\subsection{The exogenous contextual supports}
\begin{definition}[Contextual supports]\label{def:X}
Let $X_t=(S_t,C_t,K_t)$ be team-level antecedents:
$S_t$ supportive leader behavior, $C_t$ organizational context support
(resources, information, rewards aligned with learning), and $K_t$
coaching. These are the ``contextual supports'' of the problem and live at
the team/organizational level. Write $X_t\in\mathbb{R}^3$ and let
$\beta=(\beta_S,\beta_C,\beta_K)$ be a coefficient vector.
\end{definition}

\subsection{Individual-level antecedents and performance}
\begin{definition}\label{def:Z}
Let $Z_{t,i}$ collect individual-level antecedents of member $i$
(e.g.\ tenure, self-efficacy, status). Let $P_t$ denote a team
performance outcome (e.g.\ supervisor-rated effectiveness), scaled so that
larger is better.
\end{definition}

\section{The structural model}

We posit a recursive (acyclic) linear structural system. All disturbance
terms $\varepsilon_\bullet$ are mean-zero and, by the recursive ordering
$X_t\to \psi_{t,i}\to L_t\to P_t$, are uncorrelated with the regressors
that precede them in the ordering (Assumption~\ref{ass:exog}).

\begin{assumption}[Risk-cost formation of beliefs]\label{ass:safety}
For each member,
\begin{equation}\label{eq:psi}
\psi_{t,i} \;=\; \alpha_0 + \beta^{\top} X_t + \gamma^{\top} Z_{t,i}
             + \varepsilon^{\psi}_{t,i},
\end{equation}
with $\beta_S,\beta_C,\beta_K>0$. That is, each contextual support
raises members' beliefs that interpersonal risk taking is safe.
\end{assumption}

\begin{assumption}[Shared climate]\label{ass:share}
Members of a team share a common response to team-level inputs: the slopes
$\beta,\gamma$ in \eqref{eq:psi} are common across $i\in M_t$, and the
idiosyncratic shocks $\varepsilon^{\psi}_{t,i}$ are i.i.d.\ across $i$
within a team, independent of $(X_t,\bar Z_t)$, with mean $0$ and variance
$\sigma^2_\psi$. Write the team's \emph{systematic component}
\[
\tau_t \;=\; \alpha_0+\beta^\top X_t+\gamma^\top\bar Z_t,
\qquad
\bar Z_t=\tfrac1{n_t}\textstyle\sum_{i\in M_t} Z_{t,i},
\]
and let $\tau^2=\operatorname{Var}_{t}(\tau_t)$ denote its variance across
teams in the population $\mathcal T$. Two consequences follow.
\emph{(a) Within-team dispersion is bounded:} since members of a team share
the common $\tau_t$, the only source of within-team variation in
$\psi_{t,i}=\tau_t+\gamma^\top(Z_{t,i}-\bar Z_t)+\varepsilon^{\psi}_{t,i}$
is idiosyncratic, so $\Psi_t$ is a meaningful shared-climate aggregate
whenever $\tau^2$ is large relative to $\sigma^2_\psi$.
\emph{(b) Aggregation is reliable:} averaging \eqref{eq:psi} over $i$,
$\Psi_t=\tau_t+\tfrac1{n_t}\sum_i\varepsilon^{\psi}_{t,i}$, so the
reliability of the team-mean climate (the squared correlation between
$\Psi_t$ and its systematic target $\tau_t$) is
\[
\rho \;=\;
\frac{\operatorname{Var}_t(\tau_t)}
     {\operatorname{Var}_t(\tau_t)+\sigma^2_\psi/n_t}
\;=\;\frac{\tau^2}{\tau^2+\sigma^2_\psi/n_t}\;\longrightarrow\;1
\]
as $n_t\to\infty$ or as the systematic between-team variance $\tau^2$
dominates $\sigma^2_\psi$. Thus $\Psi_t$ is a well-defined shared
team-level variable and not an artifact of averaging noise.
\end{assumption}

\begin{assumption}[Behavior responds to climate]%
\label{ass:behavior}
Aggregated learning behavior is generated by
\begin{equation}\label{eq:L}
L_t \;=\; \lambda_0 + \zeta^{\top}X_t + \delta\,\Psi_t + \theta^{\top}\bar Z_t
          + \varepsilon^{L}_t,
\end{equation}
with $\delta>0$. Here $\delta$ is the (structural) effect of the shared
climate on behavior, justified by Lemma~\ref{lem:risk}, and
$\zeta=(\zeta_S,\zeta_C,\zeta_K)$ is the \emph{residual direct} effect of
the contextual supports on learning behavior that does \emph{not} pass
through psychological safety. We do \emph{not} assume $\zeta=0$ a priori;
rather, $\zeta=0$ is a substantive, testable restriction (the
\emph{complete-mediation restriction}) whose behavioral content---that the
supports affect the interpersonally risky learning activities
$B_t^1,\dots,B_t^5$ \emph{only} by changing members' beliefs about the
consequences of risk taking---is motivated, but not proved, by
Lemmas~\ref{lem:risk}--\ref{lem:supports}.
\end{assumption}

\begin{assumption}[Performance responds to learning]\label{ass:perf}
\begin{equation}\label{eq:P}
P_t \;=\; \pi_0 + \mu\, L_t + \nu^{\top} W_t + \varepsilon^{P}_t,
\qquad \mu>0,
\end{equation}
where $W_t$ are performance covariates excluded from \eqref{eq:L}.
\end{assumption}

\begin{assumption}[Strict exogeneity / no unobserved confounding]\label{ass:exog}
The structural disturbances satisfy
\begin{itemize}
\item[(A0)] $\mathbb{E}[\varepsilon^{\psi}_{t,i}\mid X_t,Z_{t,i}]=0$;
\item[(A1)] $\mathbb{E}[\varepsilon^{L}_t\mid X_t,\Psi_t,\bar Z_t]=0$;
\item[(A2)] $\mathbb{E}[\varepsilon^{P}_t\mid L_t,W_t]=0$.
\end{itemize}
The crucial condition is (A1): the behavior disturbance
$\varepsilon^L_t$ is mean-independent of \emph{all} the regressors in
\eqref{eq:L}, including $X_t$ \emph{and} the mediator $\Psi_t$. This rules
out any unobserved factor that simultaneously raises psychological safety
and learning behavior (which would render $\Psi_t$ endogenous and the
$\Psi_t\to L_t$ link unidentified). (A1) is a genuine identification
hypothesis, not a logical consequence of the behavioral axioms; it is the
price of reading the OLS coefficient on $\Psi_t$ as the structural effect
$\delta$.
\end{assumption}

These are the model's hypotheses: Assumptions~\ref{ass:safety}--\ref{ass:perf}
are the behavioral axioms (their signs justified by the mechanism lemmas of
Section~3), while Assumption~\ref{ass:exog} is the statistical
identification condition. Everything below is derived from them, and we
indicate at each step which hypothesis is used.

\section{Justification of the structural links (mechanisms)}

\begin{lemma}[Every learning activity is interpersonally risky]\label{lem:risk}
Each coordinate $B_t^k$, $k=1,\dots,5$, of Definition~\ref{def:LB} is an
act whose execution by a member can, in the absence of safety, incur an
interpersonal cost (being seen as ignorant, incompetent, negative, or
disruptive). Hence the disposition to perform $B_t^k$ is decreasing in the
\emph{perceived} interpersonal cost of risk taking and increasing in
$\psi_{t,i}$.
\end{lemma}

\begin{proof}
We argue activity by activity. (1) Asking a question makes one's lack of
knowledge observable to others; the possible cost is the image of
ignorance. (2) Seeking feedback exposes uncertainty about one's own
performance and invites evaluation; the cost is a negative appraisal. (3)
Openly discussing an error or unexpected outcome makes a failure
attributable to identifiable members; the cost is blame. (4) Reflecting on
past performance surfaces collective and individual shortcomings for
public examination; the cost is criticism. (5) Experimenting with a new
method departs from the sanctioned routine and may fail visibly; the cost
is the image of incompetence or the charge of being disruptive. In each
case the act has the form ``expose a potential deficiency / deviation
before others,'' which is the definition of an interpersonal risk. A
rational member weighs the expected interpersonal cost against the benefit;
since $\psi_{t,i}$ measures (inversely) that expected cost, the propensity
to act is monotone increasing in $\psi_{t,i}$. As $L_t$ in
Definition~\ref{def:LB} is a positive combination of the five activities,
the team propensity for $L_t$ inherits this monotonicity. This is exactly
the qualitative content encoded by $\delta>0$ in \eqref{eq:L}.
\end{proof}

\begin{lemma}[Contextual supports raise safety]\label{lem:supports}
Supportive leadership, context support, and coaching each lower the
expected interpersonal cost of risk taking, hence raise $\psi_{t,i}$:
the signs $\beta_S,\beta_C,\beta_K>0$ in \eqref{eq:psi} are justified.
\end{lemma}

\begin{proof}
(i) Supportive leader behavior ($S_t$): a leader who is accessible,
invites input, and does not punish messengers reduces the probability that
a question, an admitted error, or a failed experiment is met with
sanction; the leader's behavior is the most salient signal members use to
estimate interpersonal cost, so increasing $S_t$ decreases that estimate,
i.e.\ raises $\psi_{t,i}$ ($\beta_S>0$). (ii) Context support ($C_t$):
resources, information access, and rewards aligned with learning make risk
taking less likely to be futile or penalized and more likely to be
rewarded, lowering its net expected cost ($\beta_C>0$). (iii) Coaching
($K_t$): coaching frames difficulties as shared learning problems, models
inquiry, and provides developmental rather than evaluative feedback, again
lowering the expected interpersonal cost ($\beta_K>0$). Each mechanism acts
on the \emph{belief about consequences of risk taking}, which is precisely
$\psi_{t,i}$; this is the channel that motivates (but does not by itself
prove) the complete-mediation restriction $\zeta=0$ of
Assumption~\ref{ass:behavior} and Theorem~\ref{thm:main}(iii). The
shared-slope condition
(Assumption~\ref{ass:share}) lets these team-level inputs move the team
\emph{mean} $\Psi_t$, giving $\partial \Psi_t/\partial X_t=\beta$.
\end{proof}

\begin{lemma}[Learning behavior raises performance]\label{lem:perf}
$\mu>0$ in \eqref{eq:P}: greater learning behavior improves team
performance.
\end{lemma}

\begin{proof}
Ongoing teams in real organizations face changing tasks, technologies, and
demands. The five activities each generate or apply knowledge that adjusts
the team's process to current conditions: questions and feedback import
missing information; discussing errors and reflecting convert past
outcomes into corrected routines; experimentation discovers superior
methods. Each adjustment weakly improves the match between the team's
process and its task environment, and at least one is strictly improving
whenever the environment is non-stationary; aggregating over the
activities gives a positive marginal effect of $L_t$ on the performance
outcome $P_t$, i.e.\ $\mu>0$. The covariates $W_t$ absorb performance
determinants unrelated to learning so that $\mu$ isolates the
learning channel (Assumption~\ref{ass:exog}).
\end{proof}

\section{The mediation theorem}

We use the standard definition of (complete) mediation in a recursive
linear system.

\begin{definition}[Direct, indirect, total effects]\label{def:effects}
In the system \eqref{eq:psi}--\eqref{eq:L}, holding individual antecedents
fixed at their team means, the effect of $X_t$ on $L_t$ decomposes as
\[
\underbrace{\frac{dL_t}{dX_t}}_{\text{total}}
=\underbrace{\frac{\partial L_t}{\partial X_t}\Big|_{\Psi_t}}_{\text{direct}}
+\underbrace{\frac{\partial L_t}{\partial \Psi_t}\,\frac{\partial \Psi_t}{\partial X_t}}_{\text{indirect}} .
\]
$\Psi_t$ \emph{completely mediates} the effect of $X_t$ on $L_t$ if the
direct effect is $0$ and the indirect effect is nonzero.
\end{definition}

\begin{theorem}[Identification of the mediation structure; complete mediation as a testable restriction]%
\label{thm:main}
Suppose Assumptions~\ref{ass:safety}--\ref{ass:exog} hold and that the
population covariance matrix of $(X_t,\Psi_t,\bar Z_t)$ is nonsingular
(no collinearity; this holds whenever $\sigma^2_\psi>0$, since then
$\Psi_t$ carries belief-shock variation orthogonal to $(X_t,\bar Z_t)$).
Then:
\begin{enumerate}
\item[\textup{(i)}] \emph{(Identification.)} The population least-squares
regression of $L_t$ on $(X_t,\Psi_t,\bar Z_t)$ recovers the structural
coefficients $(\zeta,\delta)$ of \eqref{eq:L} without bias; in particular
the coefficient on $\Psi_t$ equals $\delta>0$ and the coefficient on $X_t$
equals the residual direct effect $\zeta$. The regression of $\Psi_t$ on
$(X_t,\bar Z_t)$ recovers $\beta$ \textup{(Step~1 below)}, with
$\beta\gg0$ by Lemma~\ref{lem:supports}, and the regression of $P_t$ on
$(L_t,W_t)$ recovers $\mu>0$.
\item[\textup{(ii)}] \emph{(Effect decomposition.)} The total effect of
the contextual supports on learning behavior is
\begin{equation}\label{eq:decomp}
\frac{dL_t}{dX_t}
=\underbrace{\zeta}_{\text{direct}}
+\underbrace{\delta\,\beta}_{\text{indirect (through }\Psi_t)},
\end{equation}
with the indirect (mediated) part $\delta\beta\gg0$ always strictly
positive.
\item[\textup{(iii)}] \emph{(Complete mediation iff $\zeta=0$.)} $\Psi_t$
\emph{completely} mediates the support--learning relationship---direct
effect $0$, indirect effect $\delta\beta\neq0$---\emph{if and only if} the
complete-mediation restriction $\zeta=0$ holds. In that case
\[
\frac{dL_t}{dX_t}=\delta\beta,\qquad
\frac{dL_t}{dS_t}=\delta\beta_S>0,\quad
\frac{dL_t}{dC_t}=\delta\beta_C>0,\quad
\frac{dL_t}{dK_t}=\delta\beta_K>0,
\]
and, since $\zeta$ is identified by (i), the restriction $\zeta=0$ is
\emph{empirically testable and falsifiable}: a nonzero estimated
coefficient on $X_t$ in the regression of $L_t$ on $(X_t,\Psi_t,\bar Z_t)$
refutes complete mediation.
\item[\textup{(iv)}] \emph{(Carry-through to performance.)} The mediated
chain $X_t\to\Psi_t\to L_t\to P_t$ contributes
\begin{equation}\label{eq:perfchain}
\frac{\partial P_t}{\partial X_t}\Big|_{\text{through }\Psi_t,L_t}
=\mu\,\delta\,\beta\gg0,
\end{equation}
so the total support effect on performance is $\mu(\zeta+\delta\beta)$, of
which $\mu\delta\beta\gg0$ flows through the climate--learning chain;
under complete mediation ($\zeta=0$) the entire effect is
$\mu\delta\beta$.
\end{enumerate}
\end{theorem}

\begin{proof}
\emph{Step 1 (mediator is driven by supports; recovery of $\beta$).}
Average \eqref{eq:psi} over $i\in M_t$. By Assumption~\ref{ass:share} the
slopes $\beta,\gamma$ are common, so
\[
\Psi_t=\frac1{n_t}\sum_i\psi_{t,i}
=\alpha_0+\beta^\top X_t+\gamma^\top \bar Z_t
+\bar\varepsilon^\psi_t,\qquad
\bar\varepsilon^\psi_t:=\tfrac1{n_t}\textstyle\sum_i\varepsilon^\psi_{t,i}.
\]
By (A0), $\mathbb{E}[\bar\varepsilon^\psi_t\mid X_t,\bar Z_t]=0$, so the
population regression of $\Psi_t$ on $(X_t,\bar Z_t)$ recovers
$\beta=(\beta_S,\beta_C,\beta_K)$, which is strictly positive in each
coordinate by Lemma~\ref{lem:supports}. In particular
$\partial\Psi_t/\partial X_t=\beta$.

\emph{Step 2 (identification of $(\zeta,\delta)$; proof of (i)).}
Regress $L_t$ as given by \eqref{eq:L} on the three regressor blocks
$(X_t,\Psi_t,\bar Z_t)$. By (A1),
$\mathbb{E}[\varepsilon^L_t\mid X_t,\Psi_t,\bar Z_t]=0$, i.e.\ the
disturbance is mean-independent of \emph{all} regressors that appear in
\eqref{eq:L}; and by the nonsingularity hypothesis the regressor block has
full column rank in the population. The population least-squares projection
of a variable whose conditional mean is linear in the regressors, with
mean-zero residual orthogonal to those regressors, returns exactly the
linear coefficients. Hence the regression recovers
$(\zeta,\delta,\theta)$ without bias; in particular the coefficient on
$\Psi_t$ is $\delta>0$ (Lemma~\ref{lem:risk}) and the coefficient on
$X_t$ is $\zeta$. Likewise, by (A2) the regression of $P_t$ on
$(L_t,W_t)$ recovers $\mu>0$ (Lemma~\ref{lem:perf}). This proves~(i).
Note that the nonsingularity hypothesis is essential: if $\sigma^2_\psi=0$
then $\Psi_t=\alpha_0+\beta^\top X_t+\gamma^\top\bar Z_t$ is an exact
linear function of $(X_t,\bar Z_t)$, the regressor block is
rank-deficient, and $(\zeta,\delta)$ are \emph{not} separately identified.

\emph{Step 3 (effect decomposition; proof of (ii)).}
Substitute $\Psi_t$ from Step~1 into \eqref{eq:L}:
\begin{equation}\label{eq:reduced}
L_t=\big(\lambda_0+\delta\alpha_0\big)
+\big(\zeta+\delta\beta\big)^\top X_t
+\big(\theta+\delta\gamma\big)^\top\bar Z_t
+\delta\,\bar\varepsilon^\psi_t
+\varepsilon^L_t .
\end{equation}
Differentiating, the total effect is
$dL_t/dX_t=\zeta+\delta\beta$. By Definition~\ref{def:effects} the second
term, $\delta\,\partial\Psi_t/\partial X_t=\delta\beta$, is the indirect
(mediated) effect and the first term, $\partial L_t/\partial X_t|_{\Psi_t}
=\zeta$, is the direct effect. Since $\delta>0$ and $\beta\gg0$, the
indirect part $\delta\beta\gg0$ is strictly positive in each coordinate.
This is \eqref{eq:decomp}, proving~(ii).

\emph{Step 4 (complete mediation iff $\zeta=0$; proof of (iii)).}
By Definition~\ref{def:effects}, $\Psi_t$ \emph{completely} mediates the
relationship iff the direct effect is $0$ and the indirect effect is
nonzero. By Step~3 the indirect effect is $\delta\beta\neq0$ \emph{always}
(as $\delta>0,\beta\gg0$), and the direct effect equals $\zeta$. Therefore
complete mediation holds \emph{iff} $\zeta=0$. When $\zeta=0$,
\eqref{eq:reduced} gives $dL_t/dX_t=\delta\beta$ with the stated
strictly positive coordinatewise effects $\delta\beta_S,\delta\beta_C,
\delta\beta_K$. Because $\zeta$ is identified by Step~2, the restriction
$\zeta=0$ is testable: a nonzero estimated $X_t$-coefficient in the
regression of $L_t$ on $(X_t,\Psi_t,\bar Z_t)$ refutes complete mediation.
We emphasize that $\zeta=0$ is a substantive empirical hypothesis, not a
theorem of the behavioral axioms: Lemmas~\ref{lem:risk}--\ref{lem:supports}
establish only the signs $\delta>0$ and $\beta\gg0$ (hence the positivity
and existence of the mediated path) and \emph{motivate}---but do not
prove---the absence of a residual direct path. This proves~(iii).

\emph{Step 5 (carry-through to performance; proof of (iv)).}
Substitute \eqref{eq:reduced} into \eqref{eq:P}:
\[
P_t=\big(\pi_0+\mu\lambda_0+\mu\delta\alpha_0\big)
+\mu\big(\zeta+\delta\beta\big)^\top X_t
+\mu(\theta+\delta\gamma)^\top\bar Z_t
+\nu^\top W_t+ u_t,
\]
with $u_t=\mu\delta\,\bar\varepsilon^\psi_t+\mu\varepsilon^L_t
+\varepsilon^P_t$. Hence the total support effect on performance is
$dP_t/dX_t=\mu(\zeta+\delta\beta)$, of which the portion routed through the
climate--learning chain $X_t\to\Psi_t\to L_t\to P_t$ is $\mu\delta\beta$,
strictly positive in each coordinate because $\mu>0$
(Lemma~\ref{lem:perf}), $\delta>0$, and $\beta\gg0$. This is
\eqref{eq:perfchain}. Under complete mediation ($\zeta=0$) the entire
support effect on performance is $\mu\delta\beta\gg0$, flowing solely
through that chain. This proves~(iv).
\end{proof}

\begin{corollary}[Classification of mediation regimes]\label{cor:partial}
Under the hypotheses of Theorem~\ref{thm:main}, exactly one of the
following holds, and which one is empirically determined by the identified
value of $\zeta$:
\begin{itemize}
\item \emph{Complete mediation} if $\zeta=0$: direct effect $0$, indirect
effect $\delta\beta\gg0$ \textup{(Theorem~\ref{thm:main}(iii))}.
\item \emph{Partial mediation} if $\zeta\neq0$: both the direct part
$\zeta$ and the indirect part $\delta\beta\gg0$ are nonzero, and
$\Psi_t$ accounts for the fraction
$\delta\beta/(\zeta+\delta\beta)$ of the total effect (computed
coordinatewise, where defined).
\end{itemize}
In both regimes the mediated portion $\delta\beta$ is strictly positive in
each coordinate by Lemmas~\ref{lem:risk}--\ref{lem:supports}. Thus the
model nests both regimes; complete mediation is the testable special case
$\zeta=0$ and is \emph{not} presumed.
\end{corollary}

\begin{proof}
Immediate from the decomposition $dL_t/dX_t=\zeta+\delta\beta$
\eqref{eq:decomp} and Definition~\ref{def:effects}: the indirect part
$\delta\beta\neq0$ always, so the regime is determined solely by whether
the direct part $\zeta$ vanishes. The mediated fraction is the ratio of
the indirect part to the total. Identifiability of $\zeta$ is
Theorem~\ref{thm:main}(i).
\end{proof}

\section{Multilevel summary of antecedents and mechanisms}

The model answers each part of the problem explicitly.

\begin{itemize}
\item \textbf{Individual level.} The antecedents $Z_{t,i}$ enter
\eqref{eq:psi} and, via their team mean $\bar Z_t$, \eqref{eq:L}.
Mechanism: a member's tenure, self-efficacy, and status calibrate the
\emph{personal} expected cost of interpersonal risk
($\gamma^\top Z_{t,i}$ in \eqref{eq:psi}) and the personal capacity to act
($\theta^\top\bar Z_t$ in \eqref{eq:L}). They modulate but do not
substitute for the shared climate.

\item \textbf{Team level.} Supportive leadership $S_t$ and coaching $K_t$
are team-level inputs whose mechanism (Lemma~\ref{lem:supports}) is to
lower the shared expected cost of risk taking, raising $\Psi_t$; and the
shared climate $\Psi_t$ itself is the team-level mediator whose mechanism
(Lemma~\ref{lem:risk}) is to make the interpersonally risky learning
activities $B_t^1,\dots,B_t^5$ enactable.

\item \textbf{Organizational level.} Context support $C_t$
(resources, information, learning-aligned rewards) is the
organizational-level input; its mechanism (Lemma~\ref{lem:supports}) is to
reduce the futility and penalty of risk taking, again operating through
$\Psi_t$.

\item \textbf{Mediation.} Theorem~\ref{thm:main} identifies the
mediation structure: the total support--learning effect decomposes as
$\zeta+\delta\beta$, with the mediated part $\delta\beta\gg0$ always
present and the direct part $\zeta$ identified from data. $\Psi_t$
completely mediates the relationship \emph{iff} $\zeta=0$
(Theorem~\ref{thm:main}(iii)), a testable restriction whose behavioral
content is motivated by Lemmas~\ref{lem:risk}--\ref{lem:supports}; the
partial-mediation regime ($\zeta\neq0$) is the general case
(Corollary~\ref{cor:partial}).

\item \textbf{Performance.} Equation~\eqref{eq:P} and
Lemma~\ref{lem:perf} establish $\Psi_t\to L_t\to P_t$ with total support
effect $\mu\delta\beta\gg0$ \eqref{eq:perfchain}: the climate variable
improves performance precisely by enabling learning behavior, not
independently of it.
\end{itemize}

\begin{remark}[Consistency check of the identified effects under (A0)--(A2)]
The reduced-form identities of Theorem~\ref{thm:main} are algebraic
consequences of the structural model that, under the exogeneity conditions
(A0)--(A2), become population regression restrictions. We stress that a
simulation cannot \emph{prove} complete mediation: any Monte-Carlo draw
merely reflects the value of $\zeta$ that was put into the data-generating
process. As a consistency check on the algebra of
Theorem~\ref{thm:main}(i)--(ii), a draw from the system
\eqref{eq:psi}--\eqref{eq:P} with $n_t=6$, $\delta=1.2$,
$\beta=(0.8,0.5,0.3)$, $\mu=0.9$ \emph{and the direct effect set to}
$\zeta=0$ reproduces the predicted coefficients: the regression of $L_t$
on $(X_t,\Psi_t,\bar Z_t)$ returns coefficient $\approx 0$ on each
component of $X_t$ (consistent with $\zeta=0$) and $\approx\delta$ on
$\Psi_t$; the reduced regression of $L_t$ on $(X_t,\bar Z_t)$ returns
$\approx\delta\beta=(0.96,0.60,0.36)$; and the regression of $P_t$ on
$(X_t,\bar Z_t,W_t)$ returns $\approx\mu\delta\beta=(0.864,0.54,0.324)$.
A complementary draw with $\zeta\neq0$ instead returns the nonzero direct
coefficient $\zeta$ in the first regression, illustrating that the
complete-mediation restriction $\zeta=0$ is genuinely \emph{testable and
falsifiable}, not assumed. In short, the simulation confirms the
identification algebra; the empirical truth of $\zeta=0$ in any field
sample is a separate question to be settled by estimating $\zeta$.
\end{remark}

This completes the construction and justification of the model.
\end{document}
